\chapter{Validation against goldfish}\label{ch:validation}

\nwdynam{} is a reimplementation, and the only reason to trust its numbers is that they agree with the program that it reimplements. This chapter describes how the comparison with \goldfish{}~\gfver{} (R~4.6) was done, shows a comparison that was run for this manual, and lists the places where the two programs differ. Appendix~\ref{app:translation} gives the translation of the commands and Chapter~\ref{ch:limits} lists what \nwdynam{} does not do.

\section{How the comparison was done}\label{sec:howvalidated}

The comparisons are of three kinds. The package's tests (the do-files in the directory \texttt{cscripts}) rerun the first two without R; the scripts that produced the reference values from \goldfish{} are in the directory \texttt{dev}.

\begin{enumerate}[nosep]
  \item \textbf{The compiled plugin against reference evaluators.} Every likelihood is also written in Mata, as plain reference code that does not belong to the package's estimation path and was verified against \goldfish{} while it was developed. The plugin is compared with these evaluators on random small event sequences: with simultaneous events, event values (also zero and negative), the option \opt{sequential}, windows, weighted effects, two-mode networks, \opt{opportunities()}, \opt{tie()} and every subset of the effects of a sub-model. The log likelihood and the gradient must agree to a relative $10^{-9}$ and the Hessian with the numerical derivative of the reference gradient to $10^{-5}$.
  \item \textbf{Random event logs against goldfish.} Event logs (random or hand-built), including attribute changes, memberships, windows and networks, were fitted by \goldfish{}; its estimates, standard errors and log likelihoods are stored with the package, and the tests compare \nwdynam{} with them. Table~\ref{tab:logs} lists the sets.
  \item \textbf{The emails.} The first 631 emails of the manufacturing data were fitted by both programs. Section~\ref{sec:cmpmanual} shows this comparison, rerun for this manual.
\end{enumerate}

\begin{table}[ht]
\centering\small
\begin{tabular}{@{}>{\raggedright\arraybackslash}p{6.4cm}>{\raggedright\arraybackslash}p{3.0cm}>{\raggedright\arraybackslash}p{4.8cm}@{}}
\toprule
Features & Logs & Tolerance\\
\midrule
composition change (\opt{present()}, \opt{compchange()}), together with time-varying covariates, in all sub-models, one- and two-mode & 67 random and fixed & estimates, standard errors, log likelihood $10^{-6}$\\
time-varying covariates (\opt{change()}), every attribute effect, replace and increment, changes at the time of an event, a window, several variables & 42 random & the same; coordination: estimates $10^{-5}$, standard errors 10\%\\
windowed rate effects with \opt{intercept}, \opt{tertiuswindow()}, windows and weights in the coordination sub-model, effects in two-mode networks, \opt{opportunities()} with covariates, \opt{change()} in two-mode networks and with \opt{opportunities()} & 306 random & the same; coordination: estimates $5\cdot10^{-4}$, standard errors 10\%\\
effects of other networks (\opt{tie()}, closures, mixed closures, \texttt{tertiusdiff}, \opt{ignorerep}, with windows and weights, static and changing networks) & about 300 random & the same; coordination: estimates $5\cdot10^{-4}$, standard errors 10\%\\
postestimation (\texttt{estat}): per-event log likelihoods and probabilities, \texttt{logLik}, AIC, BIC, outliers, change points & 90 random & $10^{-10}$ (values agree to $10^{-13}$); outlier indices and change points exact\\
\bottomrule
\end{tabular}
\caption{The comparisons with \goldfish{}~\gfver{} on random and hand-built event logs.}\label{tab:logs}
\end{table}

For the coordination sub-model the standard errors are numerical (from a numerical derivative of the likelihood, which is flat and not concave) in both programs and agree to 10 per cent. Effects that \goldfish{} cannot compute, or computes with an ordering of events that \nwdynam{} does not reproduce, are refused with a message (Chapter~\ref{ch:limits}), not silently approximated. The tolerances are those of the tests; a relative tolerance of $10^{-6}$ is not a claim of agreement to the last digit of every number, and a fit that ends in a different local maximum of the coordination likelihood would show up as a larger difference, not as a smaller one.

\section{A comparison run for this manual}\label{sec:cmpmanual}

To show what these comparisons look like, five models were fitted by both programs on the first 631 emails among 154 employees. The Stata side is the do-file \texttt{manual/ex/\allowbreak ex91\_compare\_stata.do}, the R side \texttt{manual/ex/r/\allowbreak compare.R}; both read the same event files, exported by \texttt{ex90\_export\_for\_r.do} (the time stamps in full precision). The models are the choice model with \opt{inertia}, \opt{recip}, \opt{indeg} (with \opt{trans} and with a window of half a day on \opt{inertia} for the emails), and the rate model with \opt{indeg} and \opt{outdeg} without and with an intercept. Table~\ref{tab:cmp} shows the estimates, the relative differences and the standard errors.

The largest relative difference of an estimate is \cmpMaxB{} (\cmpMaxWhere{}), that of a standard error \cmpMaxSE{} and that of a log likelihood \cmpMaxLL{}. The agreement does not depend on the kind of model: all differences are of the same small size in the choice and in the rate models, with and without an intercept, with a window and with a closure effect.

\begin{table}[ht]
\centering\small
{\setlength{\tabcolsep}{3pt}\begin{tabular}{@{}lrrrrr@{}}
\toprule
Effect & nwdynam & goldfish & rel.\ diff. & s.e.\ nwdynam & s.e.\ goldfish \\
\midrule
\multicolumn{6}{@{}p{14.6cm}@{}}{\textbf{Social Evolution, choice: inertia, recip, indeg}; log likelihood -695.4608 (nwdynam), -695.4608 (goldfish)}\\
\quad \texttt{inertia} & 5.36556 & 5.36556 & 4.1e-07 & 0.19375 & 0.19375 \\
\quad \texttt{recip} & 1.38944 & 1.38944 & 4.9e-07 & 0.16729 & 0.16729 \\
\quad \texttt{indeg} & -0.17353 & -0.17353 & 1.3e-06 & 0.09097 & 0.09097 \\
\midrule
\multicolumn{6}{@{}p{14.6cm}@{}}{\textbf{Social Evolution, rate: indeg, outdeg}; log likelihood -1750.9123 (nwdynam), -1750.9123 (goldfish)}\\
\quad \texttt{indeg} & 0.55144 & 0.55144 & 3.4e-08 & 0.06634 & 0.06634 \\
\quad \texttt{outdeg} & 0.26378 & 0.26378 & 9.6e-09 & 0.02839 & 0.02839 \\
\midrule
\multicolumn{6}{@{}p{14.6cm}@{}}{\textbf{Social Evolution, rate with intercept: indeg, outdeg}; log likelihood -6008.9834 (nwdynam), -6008.9834 (goldfish)}\\
\quad \texttt{Intercept} & -14.15642 & -14.15642 & 4.3e-08 & 0.06749 & 0.06749 \\
\quad \texttt{indeg} & 0.76226 & 0.76226 & 3.2e-07 & 0.06018 & 0.06018 \\
\quad \texttt{outdeg} & 0.25684 & 0.25684 & 7.8e-08 & 0.02909 & 0.02909 \\
\midrule
\multicolumn{6}{@{}p{14.6cm}@{}}{\textbf{Emails, choice: inertia, recip, indeg}; log likelihood -2307.3971 (nwdynam), -2307.3971 (goldfish)}\\
\quad \texttt{inertia} & 1.76485 & 1.76485 & 6.5e-09 & 0.14227 & 0.14227 \\
\quad \texttt{recip} & 2.92238 & 2.92238 & 9.9e-08 & 0.14011 & 0.14011 \\
\quad \texttt{indeg} & 0.10962 & 0.10962 & 4.2e-08 & 0.00831 & 0.00831 \\
\midrule
\multicolumn{6}{@{}p{14.6cm}@{}}{\textbf{Emails, choice: inertia, recip, indeg, trans}; log likelihood -2296.4872 (nwdynam), -2296.4872 (goldfish)}\\
\quad \texttt{inertia} & 1.73244 & 1.73244 & 9.3e-08 & 0.14179 & 0.14179 \\
\quad \texttt{recip} & 2.97843 & 2.97843 & 3.2e-07 & 0.13848 & 0.13848 \\
\quad \texttt{indeg} & 0.09478 & 0.09478 & 1.9e-07 & 0.00918 & 0.00918 \\
\quad \texttt{trans} & 0.50277 & 0.50277 & 5.6e-07 & 0.10322 & 0.10322 \\
\midrule
\multicolumn{6}{@{}p{14.6cm}@{}}{\textbf{Emails, choice: inertia (window 0.5), recip, indeg}; log likelihood -2312.8474 (nwdynam), -2312.8474 (goldfish)}\\
\quad \texttt{inertia} & 1.65198 & 1.65198 & 5.3e-08 & 0.13487 & 0.13487 \\
\quad \texttt{recip} & 3.22530 & 3.22530 & 5.1e-08 & 0.12690 & 0.12690 \\
\quad \texttt{indeg} & 0.12415 & 0.12415 & 4.3e-08 & 0.00779 & 0.00779 \\
\midrule
\multicolumn{6}{@{}p{14.6cm}@{}}{\textbf{Emails, rate: indeg, outdeg}; log likelihood -2986.2823 (nwdynam), -2986.2823 (goldfish)}\\
\quad \texttt{indeg} & 0.03246 & 0.03246 & 2.7e-08 & 0.01211 & 0.01211 \\
\quad \texttt{outdeg} & 0.29812 & 0.29812 & 1.9e-09 & 0.01763 & 0.01763 \\
\midrule
\multicolumn{6}{@{}p{14.6cm}@{}}{\textbf{Emails, rate with intercept: indeg, outdeg}; log likelihood -1.4359 (nwdynam), -1.4359 (goldfish)}\\
\quad \texttt{Intercept} & 0.39284 & 0.39284 & 4.0e-10 & 0.05664 & 0.05664 \\
\quad \texttt{indeg} & 0.01373 & 0.01373 & 2.1e-10 & 0.01199 & 0.01199 \\
\quad \texttt{outdeg} & 0.18574 & 0.18574 & 1.2e-10 & 0.01582 & 0.01582 \\
\bottomrule
\end{tabular}
}
\caption{Estimates and standard errors of \nwdynam{} (Stata~17, macOS) and \goldfish{}~\gfver{} (R~4.6.0) on the same event logs.}\label{tab:cmp}
\end{table}

\section{Where the two programs differ}\label{sec:differs}

\begin{itemize}[nosep]
  \item \textbf{Specification.} Effects are options in \nwdynam{}, a formula in \goldfish{}. Without effect options \nwdynam{} fits the default effects of the sub-model (Section~\ref{sec:defaults}).
  \item \textbf{Starting values and optimizer.} Both maximize the same likelihood, so that the estimates agree to the tolerance of their convergence criteria; the optimizers are different (Section~\ref{sec:technewton}).
  \item \textbf{The \texttt{tertius} effect for an actor with no in-neighbours.} The documentation of \goldfish{} describes an imputation of the average over the actors that have in-neighbours; \goldfish{}~\gfver{} itself computes 0 at estimation time, which is what \nwdynam{} does.
  \item \textbf{Simultaneous events.} \nwdynam{} follows \goldfish{} (Section~\ref{sec:ties}); the option \opt{sequential} is an addition.
  \item \textbf{Defaults.} \nwdynam{} has no formula that lists all effects; compare the translation table in Appendix~\ref{app:translation}.
  \item \textbf{Postestimation.} Nothing that \goldfish{} does not have: no simulation and no goodness of fit.
\end{itemize}
Before using \nwdynam{} for an important result, run the same model in \goldfish{}.
